\documentclass[conference]{IEEEtran}
\IEEEoverridecommandlockouts
\usepackage{cite}
\usepackage{amsmath,amssymb,amsfonts}
\usepackage{algorithm}
\usepackage{algpseudocode}
\usepackage{graphicx}
\usepackage{textcomp}
\usepackage{xcolor}
\usepackage{booktabs}
\usepackage{url}

\newtheorem{definition}{Definition}

\begin{document}

\title{Supplementary Material for\\Recall-Constrained Influence-Aware User Retrieval for Product Recommendation}
\maketitle

% ---------------------------------------------------------------------
\section*{S1\quad Preliminaries and Framework Details}
\label{sec:s1}

\textbf{BPR-MF.} BPR-MF learns user and item embeddings from implicit feedback
by maximising the pairwise ranking objective
\begin{equation}
\max_{\Theta}\sum_{(u,i,j)\in D}\ln\sigma\big(\hat{x}_{ui}-\hat{x}_{uj}\big)-\lambda\|\Theta\|^2,
\qquad \hat{x}_{ui}=\langle p_u,q_i\rangle,
\end{equation}
where $(u,i)$ is an observed interaction and $j$ a sampled negative. After row
normalisation the inner-product ranking coincides with the ranking induced by
$-\|p_u-q_i\|^2$, so an L2 index is the correct retrieval structure for these
embeddings.

\textbf{HNSW.} HNSW organises vectors into a hierarchy
of proximity graphs. A query descends greedily through the upper layers and then
runs a breadth-bounded search at level $0$, controlled by $\mathit{ef}$. The
returned list is approximate and its recall increases monotonically with
$\mathit{ef}$, which is why a high-$\mathit{ef}$ result serves as a proxy for
exact retrieval in Section IV-B of the main text.

\textbf{Influence maximization.} Given $G$ and a diffusion model, the task is to
select $K$ seeds maximising the expected number of reached nodes. The objective
is monotone and submodular, so greedy selection attains a $(1-1/e)$
approximation, and CELF accelerates it by lazy evaluation.
Our objective differs in two respects: coverage is counted over one hop, and the
seed set must satisfy a retrieval constraint.

\textbf{Offline and online stages.} Offline, the framework reuses the existing
HNSW index, stores the social graph in compressed sparse row form, and
materialises a degree-ordered user list. Online, the candidate-pool module emits
a pool that already contains influential users, the replacement module exchanges
eligible anchors of the similarity Top-$K$, and the balance module bounds every
exchange by the recall floor.

% ---------------------------------------------------------------------
\section*{S2\quad Complete Notation Table}
\label{sec:s2}
Table~\ref{tab:s2} lists every symbol used in the main text and in this
supplement, grouped by the part of the framework it belongs to.

\begin{table}[h]
\caption{Complete notation, grouped by the part of the framework it belongs to.}
\label{tab:s2}
\begin{center}
\small
\setlength{\tabcolsep}{4pt}
\begin{tabular}{|c|l|}
\hline
\textbf{Symbol} & \textbf{Description} \\
\hline
\multicolumn{2}{|l|}{\textit{Sets and graph}} \\
\hline
$U$, $I$ & user set and item set \\
$V$, $E$ & node set and edge set of the social graph \\
$G=(V,E)$ & undirected social graph \\
$N(v)$ & one-hop neighbour set of user $v$ \\
$\bar{d}$ & average degree of $G$ \\
\hline
\multicolumn{2}{|l|}{\textit{Representation and retrieval}} \\
\hline
$p_u$, $q_i$ & embedding of user $u$ and item $i$ \\
$d$ & embedding dimension \\
$q$ & query item vector \\
$d(u,q)$ & L2 distance between user $u$ and query $q$ \\
$\mathit{ef}$ & HNSW search breadth \\
$M$, $M_0$ & HNSW connectivity and its level-0 value \\
$NN_K(q)$ & exact Top-$K$ nearest neighbours of $q$ \\
$K$ & number of seed users to return \\
\hline
\multicolumn{2}{|l|}{\textit{Objective and constraint}} \\
\hline
$S$ & returned seed user set, $|S|=K$ \\
$f(S)$ & one-hop coverage, $|\bigcup_{v\in S}N(v)|$ \\
$\text{Recall@}K(S)$ & $|S\cap NN_K(q)|/K$ \\
$\tau$ & recall floor specified by the campaign \\
$S_{\text{proxy}}$ & proxy truth set obtained by ANN search \\
$\text{min\_overlap}$ & required overlap with $S_{\text{proxy}}$, $\lceil K\tau\rceil$ \\
$\text{slack}$ & $\text{current\_overlap}-\text{min\_overlap}$ \\
$I(\cdot)$ & indicator function \\
\hline
\multicolumn{2}{|l|}{\textit{Scoring and replacement}} \\
\hline
$\alpha$, $A$ & dynamic balancing weight and its maximum \\
$P$ & retrieval progress, $|C|/\mathit{ef}$ \\
$C$ & candidate set collected during retrieval \\
$\mathit{inf}(v,C)$ & marginal coverage gain of $v$ over $C$ \\
$\epsilon$ & distance slack of the hard threshold $\lambda$ \\
$\lambda$ & hard distance threshold, $D_{\text{bound}}(1+\epsilon)$ \\
$\theta$ & soft threshold of the replacement rule \\
$\beta$ & scale factor of $\theta$ \\
\hline
\end{tabular}
\end{center}
\end{table}

% ---------------------------------------------------------------------
\section*{S3\quad Candidate-Pool Retrieval}
\label{sec:s3}

\textbf{Why static degree is insufficient.} Degree ignores coverage overlap, so
high-degree users that share followers are all scored highly; it does not change
as the retrieval proceeds, so a node can keep attracting budget after its
neighbours are covered; and it is unrelated to matching quality, so admitting it
early pulls the search path out of the high-matching region. The marginal
coverage gain, evaluated against the collected set $C$,
\begin{equation}
\mathit{inf}(v,C)=\Big|N(v)\setminus\bigcup_{u\in C}N(u)\Big|,
\end{equation}
addresses all three.

\textbf{Dynamic $\alpha$.} The first design switched $\alpha$ from $0$ to $A$
once the pool was full, which destabilises the search path. Making $\alpha$ a
function of the progress $P=|C|/\mathit{ef}$ removes the jump. Four schedules
were compared: step, linear $\alpha=A\cdot P$, quadratic $\alpha=A\cdot P^{2}$,
and a threshold schedule that stays at zero for the first half of the progress.
Table~\ref{tab:s3-alpha} reports the comparison.

\begin{table}[h]
\caption{Different $\alpha(P)$ schedules (SIFT1M, $A{=}100$, $K{=}10$,
$\mathit{ef}{=}100$).}
\label{tab:s3-alpha}
\begin{center}
\begin{tabular}{|c|c|c|c|}
\hline
\textbf{Schedule} & \textbf{Recall@10} & \textbf{Avg. Inf.} & \textbf{Lat. (ms)} \\
\hline
Step      & 0.829 & 394.24 & 3.026 \\
Linear    & 0.848 & \textbf{403.09} & 2.289 \\
Quadratic & \textbf{0.849} & 395.42 & 2.202 \\
Threshold & 0.845 & 395.66 & 2.194 \\
\hline
\end{tabular}
\end{center}
\end{table}

All three continuous schedules outperform the step schedule, and the linear
schedule gives the highest coverage while keeping recall at 0.848 and latency at
2.289\,ms. It is therefore used throughout: early in the traversal $\alpha$ is
near zero and distance dominates, which preserves the navigation path; late in
the traversal $\alpha$ reaches $A$ and coverage decides, which enlarges the
influence potential of the pool.

\textbf{Amplitude.} Table~\ref{tab:s3-A} varies the amplitude of the linear
schedule.

\begin{table}[h]
\caption{Different $A$ values (SIFT1M, $K{=}10$, $\mathit{ef}{=}100$).}
\label{tab:s3-A}
\begin{center}
\begin{tabular}{|c|c|c|c|}
\hline
\textbf{$A$} & \textbf{Recall@10} & \textbf{Avg. Inf.} & \textbf{Lat. (ms)} \\
\hline
0   & \textbf{0.987} & 292.62 & 2.159 \\
10  & 0.971 & 294.87 & 1.997 \\
50  & 0.920 & 322.06 & 2.046 \\
100 & 0.832 & 396.42 & 2.231 \\
200 & 0.575 & 677.95 & 2.346 \\
300 & 0.339 & \textbf{945.21} & 2.464 \\
\hline
\end{tabular}
\end{center}
\end{table}

Recall decreases and coverage increases monotonically with $A$ while latency
stays between 2.0 and 2.5\,ms, so the amplitude moves the operating point
without changing the cost of the traversal. It is the quantity that the
query-level adaptive mechanism adjusts at run time.

\textbf{Workflow.} The upper layers are traversed by pure physical greedy
navigation. At level $0$ the entry node initialises the visited array, a
candidate min-heap keyed by hybrid score, a result max-heap, and a hash table of
covered neighbours. Each iteration pops the lowest-score node, computes $P$ and
$\alpha$, and evaluates the marginal gain of every unvisited neighbour; a
neighbour is admitted only if the result heap is not full or its score beats the
current worst. Admission increments the neighbour counters and eviction
decrements them, so the coverage state always reflects the current set. The
search stops when the candidate heap is empty, or when the best candidate is
worse than the worst result while the result heap is full.

% ---------------------------------------------------------------------
\section*{S4\quad Hierarchical Iterative Replacement}
\label{sec:s4}

\textbf{\texttt{1v1} rule.} The initial set $S$ is the HNSW Top-$K$ and
$D_{\text{bound}}$ is its largest distance. A candidate is admissible only if
$d_C\le\lambda$ with $\lambda=D_{\text{bound}}(1+\epsilon)$. For an admissible
candidate and a seed $V\in S$,
\begin{equation}
\text{Gain}(C)=\Big|N(C)\setminus\bigcup_{u\in S}N(u)\Big|,
\end{equation}
\begin{equation}
\text{Loss}(V)=\Big|\big\{w\in N(V):w\notin\bigcup_{u\in S\setminus\{V\}}N(u)\big\}\Big|,
\end{equation}
and the exchange is accepted when the candidate is no farther and adds coverage,
or when
\begin{equation}
\frac{\text{Gain}(C)-\text{Loss}(V)}{d_C-d_V}>\theta,
\end{equation}
\begin{equation}
\theta=\beta\cdot\frac{\mathrm{avg}(\text{Degree})}{\mathrm{avg}(\text{Dist})}.
\end{equation}

\textbf{Oversampling.} $3K$ candidates are fetched; the first $K$ form the
initial set and the remaining $2K$ enter the queue after the same filter
$\lambda$, which gives the stage access to users outside the physical
neighbourhood.

\textbf{\texttt{2v2} fallback.} When no single swap is accepted, the two seeds
with the smallest marginal loss are removed and replaced by the candidate with
the largest marginal gain together with the unselected candidate closest to the
query. The acceptance test is unchanged.

\begin{table}[h]
\caption{Effect of $\beta$ on the three levels (SIFT1M, $K{=}10$,
$\epsilon{=}0.2$).}
\label{tab:s4-beta}
\begin{center}
\scriptsize
\setlength{\tabcolsep}{2.2pt}
\begin{tabular}{|c|c|c|c|c|c|c|c|c|c|}
\hline
& \multicolumn{3}{c|}{\texttt{1v1}} & \multicolumn{3}{c|}{+OS} & \multicolumn{3}{c|}{+OS+\texttt{2v2}} \\
\cline{2-10}
$\beta$ & Rec. & Inf. & Lat. & Rec. & Inf. & Lat. & Rec. & Inf. & Lat. \\
\hline
0.2 & 0.067 & 1172.16 & 60.05 & 0.067 & 1178.43 & 59.86 & 0.074 & \textbf{1206.55} & 114.78 \\
0.5 & 0.116 & 1105.05 & 43.80 & 0.112 & 1112.96 & 46.86 & 0.124 & \textbf{1157.15} & 87.40 \\
0.8 & 0.149 & 1055.01 & 36.86 & 0.136 & 1074.11 & 40.68 & 0.154 & \textbf{1104.04} & 71.28 \\
1   & 0.168 & 1014.85 & 33.48 & 0.148 & 1062.96 & 37.12 & 0.175 & \textbf{1082.50} & 63.51 \\
2   & 0.293 & 848.11  & 22.90 & 0.259 & 890.26  & 25.18 & 0.305 & \textbf{926.65}  & 42.96 \\
5   & 0.544 & 617.16  & 13.56 & 0.519 & 638.40  & 14.08 & 0.537 & \textbf{660.38}  & 25.99 \\
10  & 0.768 & 399.72  & 7.14  & 0.748 & 406.78  & 7.97  & 0.736 & \textbf{421.95}  & 14.51 \\
20  & 0.909 & 324.65  & 4.70  & 0.888 & 348.58  & 5.51  & 0.883 & \textbf{370.66}  & 10.22 \\
50  & 0.963 & 306.72  & 4.50  & 0.956 & 312.00  & 4.69  & 0.956 & \textbf{312.00}  & 8.25 \\
100 & 0.980 & 294.86  & 4.28  & 0.976 & 296.01  & 4.52  & 0.976 & \textbf{296.01}  & 8.05 \\
\hline
\end{tabular}
\end{center}
\end{table}

Coverage increases strictly with the number of levels at every $\beta$, while
recall stays within 0.02 of the \texttt{1v1} level, so the two optimisations
raise the reachable coverage rather than shifting the recall--coverage curve.
Latency grows with the number of accepted replacements. Small $\beta$ makes the
soft threshold permissive and recall falls to 0.07; $\beta{=}100$ leaves almost
no room for replacement and the result approaches the plain ANN baseline. The
main text uses $\beta{=}20$; in the complete framework the permitted amount of
replacement is derived from the remaining slack instead.

% ---------------------------------------------------------------------
\section*{S5\quad End-to-End Procedure}
\label{sec:s5}

Algorithm~\ref{alg:s5} expands Algorithm~1 of the main text: it makes the state
shared by the three modules explicit and states the termination condition of the
replacement loop. The invariant that follows explains why the recall floor
cannot be violated.
\algrenewcommand\algorithmicrequire{\textbf{Require:}}
\begin{algorithm*}[t]
\caption{Recall-constrained influence-aware retrieval (expanded)}
\label{alg:s5}
\begin{algorithmic}[1]
\Require $q$, $K$, $\tau$, $\mathit{ef}$, $A$
\Ensure $S$ with $|S|=K$ and $|S\cap S_{\text{proxy}}|\ge\lceil K\tau\rceil$
\State $S_{\text{proxy}}\leftarrow\textsc{Ann}(q,K,\mathit{ef})$; \ $S\leftarrow S_{\text{proxy}}$
\State $m\leftarrow\lceil K\tau\rceil$; \ $\mathit{cov}\leftarrow\textsc{Counters}(S)$
\State $C\leftarrow$ hybrid-guided retrieval with dynamic $\alpha$ (S3)
\State $Q\leftarrow\{c\in C: d(c,q)\le\lambda\}$
\Repeat
  \State $v^\ast\leftarrow$ farthest eligible anchor of $S$
  \State $c^\ast\leftarrow\arg\max_{c\in Q}\text{Gain}(c)$
  \If{$\text{slack}(S)>0$ \textbf{and} $\textsc{Verify}(v^\ast,c^\ast)$}
     \State $S\leftarrow S\setminus\{v^\ast\}\cup\{c^\ast\}$; update $\mathit{cov}$
  \ElsIf{no admissible \texttt{1v1} exchange}
     \State try one \texttt{2v2} exchange; accept if coverage increases
  \Else
     \State \textbf{break}
  \EndIf
\Until{no further improvement}
\Return $S$
\end{algorithmic}
\end{algorithm*}

\textbf{Invariant.} Let $\mathit{ov}(S)=|S\cap S_{\text{proxy}}|$. Initially
$\mathit{ov}(S)=K\ge m$. An exchange is executed only when \textsc{Verify}
accepts it, and \textsc{Verify} rejects any exchange for which
\begin{equation}
\mathit{ov}'=\mathit{ov}-I(v^\ast\in S_{\text{proxy}})+I(c^\ast\in S_{\text{proxy}})
\end{equation}
falls below $m$; the \texttt{2v2} fallback exchanges non-anchor members only and
therefore leaves $\mathit{ov}(S)$ unchanged. Every reachable state satisfies the
invariant, so the returned set meets the floor by construction. The invariant is
stated against the proxy set, whose relation to exact neighbours is governed by
the recall of the ANN search; this is why the proxy is calibrated offline
against brute-force ground truth in the main text. The loop terminates because
each accepted exchange increases coverage, coverage is bounded by $|V|$, and the
number of accepted exchanges per query is bounded by $|C|$.

% ---------------------------------------------------------------------
\section*{S6\quad Datasets, Preprocessing and Baselines}
\label{sec:s6}

\textbf{SIFT.} 128-dimensional uint8 descriptors, four prefixes of 1M, 10M, 50M
and 100M vectors sharing one social graph (average degree 10) and one set of 100
queries. Squared L2 distances over uint8 are integers below $2^{24}$, so the
integer and floating-point pipelines give bit-identical results. Ground truth is
computed by brute force over the full corpus.

\textbf{CiaoDVD.} 16{,}043 users, 12{,}787 items and their trust relations; we
keep ratings above 3.5, build a leave-one-out split, and learn 64-dimensional
BPR-MF embeddings with 20 epochs, learning rate 0.05 and regularisation
$10^{-4}$. After projection onto users carrying at least one trust edge, the
graph has 2{,}310 nodes and 21{,}400 undirected edges. Queries are items,
candidates are users, and relevance is the exact $K$-nearest users of the item
embedding.

\textbf{KuaiRand-1K.} 1{,}000 users, 7{,}388 candidate videos, 40{,}726
training impressions, 10{,}316 calibration impressions, 43{,}028 random-exposure
test impressions, and a video co-click graph with 67{,}301 edges. For each
impression a matching score is computed, candidates are sorted, and CTR@10, AUC
and Logloss are measured against the real click labels.

\textbf{Baselines.} \emph{HNSW} returns the Top-$K$ of the graph index with
$\mathit{ef}{=}300$ and receives no influence signal. \emph{Degree-greedy}
runs the same replacement with the same pool, oversampled queue, $\beta$ and
floor, and differs only in replacing the lowest-degree eligible anchor instead
of the farthest one. \emph{CELF-near} maximises one-hop coverage over the
physical neighbours of the query; \emph{CELF-global} does the same over the whole
user set without a retrieval constraint.

\textbf{Execution environment.} A single node with 16 vCPUs and 60\,GB of memory.
The index is built serially (\texttt{OMP\_NUM\_THREADS=1}) so that the
construction trajectory is deterministic, and queries are issued one at a time.
The implementation is C++17, compiled with \texttt{g++ 11.2} at
\texttt{-O3 -march=native}, with $M{=}16$ and
$\mathit{ef}_{\text{construction}}{=}200$. No GPU is used.

% ---------------------------------------------------------------------
\section*{S7\quad Complete Module Ablation}
\label{sec:s7}

Table~\ref{tab:s7-a} reports the replacement module built up component by
component, and Table~\ref{tab:s7-b} compares each single module with the fused
two-stage framework.

\begin{table}[h]
\caption{Replacement module built up incrementally (SIFT1M, $K{=}10$,
$\beta{=}20$).}
\label{tab:s7-a}
\begin{center}
\begin{tabular}{|l|c|c|c|}
\hline
\textbf{Configuration} & \textbf{Recall@10} & \textbf{Avg. Inf.} & \textbf{Lat. (ms)} \\
\hline
Initial $K$ + basic \texttt{1v1} & \textbf{0.909} & 354.65 & 4.70 \\
+ oversampled pool              & 0.888 & 368.58 & 5.51 \\
+ \texttt{2v2} fallback          & 0.883 & \textbf{370.66} & 10.22 \\
\hline
\end{tabular}
\end{center}
\end{table}

\begin{table}[h]
\caption{Single modules versus the fused framework (SIFT1M, $K{=}10$,
$A{=}200$, $\beta{=}20$).}
\label{tab:s7-b}
\begin{center}
\begin{tabular}{|l|c|c|c|}
\hline
\textbf{Configuration} & \textbf{Recall@10} & \textbf{Avg. Inf.} & \textbf{Lat. (ms)} \\
\hline
Hybrid-score retrieval only & 0.594 & \textbf{688.02} & 16.83 \\
Replacement only            & 0.883 & 370.66 & 10.22 \\
Full two-stage framework    & 0.818 & 436.25 & \textbf{12.01} \\
\hline
\end{tabular}
\end{center}
\end{table}

Oversampling raises coverage by 3.9\% for a 2.3\% recall loss and 0.81\,ms of
latency, since the extra candidates are the influential users the physical
Top-$K$ lacks. The \texttt{2v2} fallback adds 0.56\% with almost no recall loss,
because it can leave states in which no single swap helps. In the fused
framework, the retrieval-only variant reaches the highest coverage (688.02) but
collapses recall to 0.594, while the replacement-only variant keeps recall at
0.883 but stalls at 370.66; the fused framework sits between them (0.818,
436.25) and has the lowest latency of the three, since the first stage already
filters the pool the second stage examines.

% ---------------------------------------------------------------------
\section*{S8\quad Influence-Maximisation Baselines}
\label{sec:s8}

\begin{table}[h]
\caption{Comparison with influence-maximisation baselines.}
\label{tab:s8}
\begin{center}
\small
\setlength{\tabcolsep}{3pt}
\begin{tabular}{|l|c|c|c|c|c|c|}
\hline
& \multicolumn{2}{c|}{CELF-near} & \multicolumn{2}{c|}{CELF-global} & \multicolumn{2}{c|}{Ours} \\
\cline{2-7}
\textbf{Setting} & Rec. & Inf. & Rec. & Inf. & Rec. & Inf. \\
\hline
SIFT 1M, $K{=}10$   & 0.490 & 240.50  & 0.000 & 1546.0 & \textbf{0.700} & \textbf{612.29} \\
SIFT 1M, $K{=}50$   & 0.499 & 1162.78 & 0.000 & 7511.0 & \textbf{0.700} & \textbf{3011.70} \\
SIFT 100M, $K{=}10$ & 0.488 & 252.27 & 0.000 & 1652.0 & \textbf{0.699} & \textbf{662.28} \\
SIFT 100M, $K{=}50$ & 0.459 & 1190.53 & 0.000 & 8219.0 & \textbf{0.698} & \textbf{3229.45} \\
CiaoDVD, $K{=}10$   & 0.630 & 47.88   & 0.000 & 720.0  & \textbf{0.700} & \textbf{413.44} \\
CiaoDVD, $K{=}20$   & 0.684 & 85.48   & 0.002 & 981.0  & \textbf{0.700} & \textbf{574.52} \\
CiaoDVD, $K{=}50$   & 0.727 & 190.00  & 0.004 & 1427.0 & \textbf{0.701} & \textbf{877.90} \\
\hline
\end{tabular}
\end{center}
\end{table}

The two variants fail in opposite directions. CELF-global reaches 8219 coverage
at $K{=}50$ on SIFT but returns a set that depends on the graph alone, so its
recall is indistinguishable from zero. CELF-near keeps recall between 0.46 and
0.73 but its coverage is an order of magnitude below ours (240.50 against
612.29 at $K{=}10$), because the physical neighbours of an item are selected for
similarity and rarely contain influential users. The proposed framework is the
only configuration that satisfies a recall floor while improving coverage over
the retrieval-only baseline.

\bibliographystyle{IEEEtran}

\end{document}
